\documentclass[10pt]{article}

\usepackage[utf8]{inputenc}

\usepackage[T1]{fontenc}

\usepackage{amsmath}

\usepackage{amsfonts}

\usepackage{amssymb}

\usepackage[version=4]{mhchem}

\usepackage{stmaryrd}



\begin{document}

Алгоритм вычислений по Э. м. легко программируется и удобен для реализации на ЭВМ.



Метод предложев Л. Эйлером (L. Euler, 1768).



JIum. [1] д е м и д о в и ч Б. П.. М а р о н И. А., Основы вычислительной математики, М., 1960. И. Б. Вапнярский.



ЭЙЛЕРА МЕТОД СУММИРОВАНИЯ - один из методов суммирования числовых и функциональных рядов. Ряд



\[

\sum_{n=0}^{\infty} a_{n}

\]



суммируем м е тод о м с ум м и ро в а н и я Э йл е р а $((E, q)$-суммируем) к сумме $S$, если



\[

\lim _{n \rightarrow \infty} \frac{1}{(q+1)^{n}} \sum_{k=0}^{n} C_{n}^{k} q^{n-k} S_{k}=S

\]



где $q>-1, S_{k}=\sum_{n=0}^{k} a_{n}$.



Впервые метод при $q=1$ применялся Л. Эїлером (L. Euler) для суммирования медленно сходящихся и расходящихся рядов. На произвольные значения $q$ метод был распространен К. Кноппом [1], поэтому Э. м. с. при любом $q$ наз. также ме т о д о м с у м м иро в а ни я Эй ле р а - К н о пा па. Э. м. с. регулярен при $q \geqslant 0$ (см. Регулярные методы суммирования); если ряд $(E, q)$-суммируем, то он суммируем и методом $\left(E, q^{\prime}\right)$ при $q^{\prime}>q>-1$ к той же сумме (см. Включение методов суммирования). При $q=0$ суммируемость Э. м. с. ряда (*) означает сходимость этого ряда. Если ряд $(E, q)$-суммируем, то его члены $a_{n}$ удовлетворяют условию $a_{n}=o\left\{(2 q+1)^{n}\right\}, q \geqslant 0$. Э. м. с. применяется также для аналитит. продолжения фуункции, определенной степенным рядом, за круг сходимости. Так, ряд $\sum_{n=0}^{\infty} z^{n}(E, q)$-суммируем к сумме $1 /(1-z)$ в круге с дентром в точке $-q$ и радиусом, равным $q+1$.



Jum.: [1] K n o p p K., "Math. Z.», 1922, Bd 15, S. 226-53; [2] е г о ж е, там же, 1923, Bd 18, S. 125-56; [3] Х а p д и Г.,

Расходящиеся ряды, пер. с англ., М., 1951; [4] Б а р о н С., Расходящиеся ряды, пер. с англ., М., 1951; [4] Б а р о н С.,



ЭИЛЕРА МНОГОЧЛЕНЫІ - многочлены вида



\[

E_{n}(x)=\sum_{k=0}^{n} C_{n}^{k} \frac{E_{k}}{2^{k}}\left(x-\frac{1}{2}\right)^{n-k},

\]



где $E_{k}$ - эйлеровь числа. Ә. м. можно последовательно вычислить по формуле



\[

E_{n}(x)+\sum_{s=0}^{n} C_{n}^{s} E_{s}(x)=2 x^{n}

\]



В частности,



\[

E_{0}(x)=1, E_{1}(x)=x-\frac{1}{2}, E_{2}(x)=x(x-1) .

\]



Э. м. удовлетворяют разностному уравнению



\[

E_{n}(x+1)+E_{n}(x)=2 x^{n}

\]



и принадлежат классу $A$ ппеля многочленов, т. е. удовлетворяют соотношению



\[

\frac{d}{d x} E_{n}(x)=n E_{n-1}(x) .

\]



Производящая функция для Э. м.:



\[

\frac{2 e^{x t}}{e^{t}+1}=\sum_{n=0}^{\infty} \frac{E_{n}(x)}{n !} t^{n} .

\]



Для Э. м. справедливо разложение в ряд Фурье



\[

\begin{gathered}

E_{n}(x)=\frac{n !}{\pi^{n+1}} \sum_{k=0}^{\infty} \frac{\cos \left[(2 k+1) \pi x+\frac{n+1}{2} \pi\right]}{(2 k+1)^{n+1}}, \\

0 \leqslant x \leqslant 1, \quad n \geqslant 1 .

\end{gathered}

\]



Э. м. удовлетворяют соотношениям:



\[

\begin{gathered}

E_{n}(1-x)=(-1)^{n} E_{n}(x), \\

E_{n}(m x)=m^{n} \sum_{k=0}^{m-1}(-1)^{k} E_{n}\left(x+\frac{k}{m}\right), \\

E_{n}(m x)=-\frac{2 m^{n}}{n+1} \sum_{k=0}^{m-1}(-1)^{k} E_{n+1}\left(x+\frac{k}{m}\right), \\

\text { если } m \text { четно. }

\end{gathered}

\]



Периодич. функции, совпадающие с правой частью (*), являются әкстремальными в Колмогорова неравенстве и ряде других әкстремальных задач теории функций. Рассматриваются также обобщенные Э. м.



лит.: [1] Э ї Ј e p J., Дифференциальное исчисление, пер. c лат., M. - J., 1949; 12$]$ Norl und N. E., Vorlesungen über ЭЙЛЕРА ПОДСТӒНОВКА - замена переменной $x=$ $=x(t)$ в интеграле



\[

\int R\left(x, \sqrt{a x^{2}+b x+c}\right) d x

\]



где $R(\cdot, \cdot)$ - рациональная функция своих аргументов, сводящая этот интеграл к интегралу от рациональной функции и имеющая один из следующих трех видов.



П е р в а я под с танов к а Эй ле ра: если $a>0$, то



\[

\sqrt{a x^{2}+b x+c}= \pm x \sqrt{a} \pm t .

\]



Вторая под с тановка Эй ле ра: если корни $x_{1}$ и $x_{2}$ квадратного трехчлена $a x^{2}+b x+c$ действительные, то



\[

\sqrt{a x^{2}+b x+c}= \pm t\left(x-x_{1}\right) .

\]



Третья под с тановка Э йле ра: если $c>0$, то



\[

\sqrt{a x^{2}+b x+c}= \pm \sqrt{c} \pm x t

\]



(в правых частях равенств можно брать любые комбинации знаков). При всех Э. п. как старая переменная интегрирования $x$, так и радикал $\sqrt{a x^{2}+b x+c}$ рационально выражаются через новую переменную $t$.



Две первые Э. п. позволяют всегда свести интеграл (1) к интегралу от рациональной функции на любом промежутке, на к-ром радикал $\sqrt{a x^{2}+b x+c}$ принимает только действительные значения.



Геометрич. смысл Э. п. состоит в том, что кривая 2-го порядка



\[

y^{2}=a x^{2}+b x+c

\]



имеет рациональное параметрич. представление: именно, если за параметр $t$ взять угловые коәфффициенты пучка секущих $y-y_{0}=t\left(x-x_{0}\right)$, проходящих через точку $\left(x_{0}, y_{0}\right)$ кривой (2), то координаты этой кривой будут рационально выражаться через переменную $t$. В случае, когда $a>0$ и, следовательно, кривая (2) является гиперболой, для того, чтобы получить 1-ю Э. Іі., за точку $\left(x_{0}, y_{0}\right)$ следует взять одну из бесконечно удаленных точек, определяемых направлениями асимптот этой гиперболы; в случае, когда корни $x_{1}$ и $x_{2}$ квадратичного трехчлена $a x^{2}+b x+c$ действительны, для того, чтобы получить 2-ю Э. п., надо взять за точку $\left(x_{0}, y_{0}\right)$ одну из точек $\left(x_{1}, 0\right)$ или $\left(x_{2}, 0\right)$; а в случае, когда $c>0$, чтобы получпть 3-ю Э. п.- одну из точек пересечения кривої (2) с осью ординат, т. е. одну из точек $(0, \pm \sqrt{c})$.



ЭЙЛЕРА ПОСТОЯННАЯ - число $а$. Д. Кудрявчев. равенством



$c=\lim _{n \rightarrow \infty}\left(1+\frac{1}{2}+\frac{1}{3}+\ldots+\frac{1}{n}-\ln n\right) \approx 0,57721566490 \ldots$,





\end{document}